\section{Method}
\label{sec:method}

\model connects a task-specific prompt modification to the evidence needed to
explain and assess it. Given natural programming requests, a generator model,
and a fixed evidence-acquisition budget, Stage I represents requirements in their
task context. Stage II constructs and prioritizes requirement hypotheses, then
estimates the effects of frozen prompt interventions on held-out tasks. Stage III
organizes the linked task semantics, edits, and effect evidence into explanations
with explicit conditions for security guidance. The same requirement scope and
non-target task requirements connect all three stages.

Consider a request to delete database records for a supplied user identifier
and return the deletion count. An explanation should identify the supplied
identifier as a SQL value, locate the parameterization requirement at the deletion
operation, and describe a modification that retains the deletion and return
behavior. It must then distinguish the reason for considering that change from
evidence about its effect. Observational prioritization supplies the former;
randomized intervention supplies the latter.

These reader-facing stages cover seven scientific operations: representation,
prioritization, hypothesis freeze, intervention/\allowbreak randomization,
measurement, outcome assembly, and inference/\allowbreak reporting.
Stage III presents the reporting
output of this same path; it adds neither an experimental role nor another
causal analysis. Development fixes the representation
and measurement rules, discovery prioritizes hypotheses, and confirmation tests
the selected policies. These roles contain disjoint deduplicated \emph{task
units} and frozen near-duplicate groups; instances and repeated generations
belonging to one unit never cross the role boundary.

\subsection{Stage I: Structured Graph Representation}

\subsubsection{Typed Symbolic Representation}
A requirement name alone does not identify an editable part of a task. In the
running example, the relevant object is the supplied identifier used as a SQL
value, the operation is deletion, and the required non-target behavior includes
returning the deletion count. Prompt TSG represents these distinctions in a
\emph{typed symbolic graph representation language}.

The language, vocabulary, and task graph have different roles. The language
defines generic node and relation types and their permitted connections. The
vocabulary supplies reusable concepts, such as SQL value parameterization.
A task graph instantiates those concepts at particular operations, objects,
and source locations. Thus, two occurrences of the same concept can refer to
different parts of a request without becoming one interchangeable graph node.
For the complete baseline input $P$ and fixed extractor $\mathcal E$, write
\begin{equation}
  G_P=T_{\mathcal E}(P)=(V,E,\tau_V,\tau_E,\alpha),
\end{equation}
where $\tau_V$ and $\tau_E$ assign node and relation types, and $\alpha$ records
concept identities, source spans, and task-local instance identities.
Nodes distinguish objects, operations, functional and safety requirements,
constraints, condition predicates, and presentation controls. Relations express input use, produced
objects, operation order, requirement scope, and conditions.

This representation keeps the task context separate from the requirement that
may be edited. In the example, adding a parameterization requirement does not
authorize changing which records are deleted or what the operation returns.
Prompt-TSG edges describe these semantic relationships; they are not causal
edges and do not establish whether the requirement affects generated code.

\subsubsection{Graph Construction and Grounding}
The graph must represent the input actually supplied to the code generator.
We therefore prepare the complete system and user messages once and construct
the baseline graph from that source. Source-only LLM annotation first identifies
the task's objects, operations, requirements, and conditions with supporting
source spans. A subsequent binding step connects these fixed local instances:
which objects an operation uses or produces, which operations and objects a
requirement governs, and which predicates qualify an action or requirement.
The binding step cannot add or relabel source facts to complete a desired
relation. Deterministic compilation checks type compatibility, instance
references, and cited evidence locations. Generated code, security outcomes,
and selector outputs do not participate in either step.

Source grounding preserves the difference between a reusable concept and its
local meaning. A parameterization requirement attached to another database
operation cannot establish that the deletion input is covered. Likewise, a
mention of an input being used is not evidence that it must be validated.
These distinctions are checked through the relevant operation and object
bindings rather than keyword co-occurrence. The same bindings let the final
explanation name the input and operation concerned, rather than presenting an
unscoped concept label.

Conditions are kept distinct from unconditional implementation constraints.
For example, an instruction to use Python does not make a database operation
conditional on Python being used. A relation must preserve the source meaning,
not merely connect endpoints of permitted types.

The vocabulary is induced from full development prompts and reviewed before
discovery; formal use fixes the vocabulary, extractor, and query rules.
Structural checks do not prove that an annotation understood the source
correctly, so source-semantic qualification is also required for the admitted
profiles. Unmodeled content and ambiguous evidence remain explicit rather
than being repaired from later outcomes.

\subsubsection{Context and Requirement Queries}
Queries turn the representation into task-scoped observations for hypothesis
construction. A context query asks whether the task contains the relevant
operation and security boundary. A factor query asks whether the source
expresses a particular requirement in that scope. The context is defined
independently of the target requirement's state: a SQL operation using a supplied
value remains the same context whether or not parameterization is requested.

After graph construction, a source-only assessment evaluates fixed operation,
input, and condition scopes against the complete source. It separately judges
whether the requirement applies in each scope and whether the source expresses
it there. Both judgments refer to existing graph instances; the assessment
cannot add facts or choose a different scope. In the SQL example, the supplied
identifier's stated role as a data value supports applicability. If an input's
role as a value or a database identifier is genuinely ambiguous, the absence of
parameterization wording cannot resolve that ambiguity. Separating applicability
from expression prevents an unexpressed but applicable requirement from being
confused with an inapplicable one.

Every query returns \textsc{present}, \textsc{absent},
\textsc{not-applicable}, or \textsc{unresolved}. For the running example,
parameterization can be absent only when its scope is applicable and the
source assessment resolves the requirement as unexpressed. A missing graph
match alone is not enough. Evidence from different operations cannot be joined
to create a positive match, and ambiguity is not silently converted to absence.
Stage I thus supplies grounded task graphs, scope bindings, and query states.
These identify where an explanation and a possible edit belong, and which
non-target requirements the edit must preserve. They do not yet establish
whether that change affects generated code.

\subsection{Stage II: Requirement-Level Causal Analysis}
\label{sec:method-causal}

\subsubsection{Hypotheses and Observational Prioritization}
A testable hypothesis connects a requirement identified in Stage I to a proposed
\emph{intervention}: a controlled change to its expression in the prompt, with
non-target task requirements held fixed. It specifies the requirement, its
scope, and whether the intervention will add an unexpressed requirement or
remove an expressed one. In the running example, the addition hypothesis asks
whether explicitly requiring parameterization for the deletion input changes
the security outcome. The factor definition fixes the scope across tasks;
local operation and object identities instantiate that definition rather than
creating a new factor for each variable name.

An \emph{Atomic} factor is one independently editable semantic requirement,
not necessarily one sentence, node, or edge. Value parameterization and
identifier allowlisting, for example, are different decisions. Requirement
addition and removal define separate hypotheses because they start from
different source states and ask different intervention questions. Both require
a concrete editing policy and a matched comparison.

A \emph{Pair} hypothesis combines two distinct requirements that admit a joint,
task-preserving comparison. It has its own applicability and compatibility
conditions. Neither requirement must first be selected, ranked, or supported
as an Atomic candidate. This keeps interactions eligible even when an Atomic
selection step would not have prioritized their individual factors. Natural
sources must establish the relevant boundary and non-target behavior, but need
not already contain experimental controls. Incomplete functional specifications
limit functionality claims without automatically invalidating a measurable
security comparison.

The analysis budget limits how many hypotheses can receive independent effect
evidence. Prioritization separates three decisions: whether natural tasks support the
comparison, whether the candidate meets the selector's structural criterion,
and how it ranks among admitted candidates. Before
using outcomes to prioritize candidates, one discoverability rule family checks
context applicability, natural-state support, source-lineage overlap, and valid
task-unit folds. Atomic candidates need support for their factor states; Pair
candidates need their own compatible four-cell support. These checks use
natural task observations, not states manufactured by experimental edits.

For Atomic hypotheses, the structural criterion comes from Fast Causal
Inference (FCI). FCI uses conditional-independence relations to infer a graph
that allows for unobserved common causes~\cite{spirtes1995fci,zheng2024causallearn}.
\model builds this observational analysis separately for each security scope
and generator model, using natural prompt-query states, approved pre-treatment
metadata, independently measured baseline outcomes, and frozen typed background
constraints. A candidate passes the structural gate when its feature--outcome
adjacency meets the qualified, frozen criterion. This graph is distinct from
the task-semantic Prompt TSG: FCI supplies selection evidence about relationships
among measured variables, not new semantic edges within a prompt.

The ranking score answers a different question: how much do outcome rates differ
between the relevant natural requirement states? A \emph{risk difference} (RD)
is a difference between outcome probabilities, expressed in percentage points.
For the primary endpoint, the outcome is a request yielding valid code that can
be assessed as secure (Section~\ref{sec:method-measurement}). The Atomic score
compares the operation-specific target state with its baseline state using
cross-fitting and covariate standardization. For addition, the comparison is
expressed minus unexpressed; for removal, it is unexpressed minus expressed.
Candidates are ranked by the absolute first-order RD, retaining the sign for
interpretation without requiring a favorable direction. RD is computed from the
discovery observations, not from FCI edge strengths; FCI determines structural
admission, while RD determines priority.

For Pair hypotheses, the same admission--ranking separation uses a different
structural source. A Prompt-TSG relation gate checks whether the two factors
have the specified semantic relationship, such as a shared sink or distinct
control points; it does not require their Atomic FCI admission. The ranking
score is an absolute cross-fitted, standardized second-order RD. It compares
the outcome-rate difference for one factor across the two states of the other,
using all four natural factor-state cells. Thus, the Pair score targets
interaction rather than the sum of two Atomic scores. Neither structural gate
nor observational RD identifies the effect of an assigned prompt change.
Generated code is evaluated to obtain outcomes, not inserted as a graph variable
or treated as a mediator.

\subsubsection{Frozen Requirement Interventions}
Atomic and Pair selectors retain budgets $K_A$ and $K_I$. The union of selected
hypotheses undergoes one common conversion into concrete prompt policies.
Each unique selected model-bound effect is confirmed once, and its result is
reused by every selector slot that refers to it. Overlap can reduce duplicate
evaluation, but cannot change a selector's denominator.

Before confirmation outcomes are available, each hypothesis fixes its scope,
intervention direction, editing recipes and weights, comparison arms, eligible
population rules, endpoint, and inference plan. Empty, failed, and
non-evaluable slots remain in the original budget without replacement.
The confirmation set cannot be reranked after seeing which changes succeed.

Natural prompts that express different requirements can also differ in other
ways. Confirmation therefore tests each frozen hypothesis by varying the
assigned prompt policy while holding the task fixed. The intervention acts on
the prompt, not directly on the code that the model will generate.

There are two intervention directions. \emph{Requirement addition} (ADD)
introduces a requirement whose absence has been resolved in the applicable
source scope. It does not assume that baseline code lacks the corresponding
protection. \emph{Requirement removal} (REMOVE) withdraws an identifiable source
requirement through a task-preserving neutral counterpart; it does not instruct
the model to implement the opposite, insecure behavior. In the SQL example,
addition explicitly requests parameterization of the supplied identifier,
whereas removal withdraws an existing parameterization requirement. Both retain
the deletion and return behavior. The two directions remain separate policies
and effect estimates, not interchangeable evidence about one symmetric effect.

For an Atomic hypothesis, the primary comparison is between a target edit $T$
and an operation-matched no-op $N$. The no-op applies a matched rewrite without
changing whether the target requirement is expressed: it leaves the requirement
unexpressed for an addition policy and retains it for a removal policy. This
comparison distinguishes the target requirement change from the matched editing
operation. The primary contrast is target minus matched no-op, not target minus
the untouched source prompt.

Two additional controls address alternative explanations: a length-matched
placebo or sham edit, and a direction-compatible generic-security instruction.
Together with $T$ and $N$, these form the Atomic four-arm design. Their purpose
is to distinguish a named requirement from rewriting, presentation changes,
or nonspecific security guidance. Specificity conclusions require their own
prospectively fixed secondary inference rules.

Pair hypotheses instead use a factorial design. For factor $j$, let $T_j$ be
its target operator and $N_j$ its matched no-op. The four cells are
\begin{equation}
\begin{aligned}
  A_{00}&=N_1+N_2, & A_{10}&=T_1+N_2,\\
  A_{01}&=N_1+T_2, & A_{11}&=T_1+T_2.
\end{aligned}
\end{equation}
This separates changing either requirement from changing both. The joint-target
cell cannot introduce a third requirement. If the operators have not been shown
to commute, the policy must cover both application orders with positive frozen
weights or exclude the pair.

A policy can have several concrete editing recipes, called \emph{realizations},
that specify wording and, where relevant, operator order. Their weights define
the policy mixture to be evaluated. Before task-specific text is materialized,
source-stratified balanced allocation assigns exactly one realization to each
task--policy coordinate. Only its complete four-arm bundle is produced and
checked. A pre-assignment failure is recorded without trying another
realization, replacing the task, or renormalizing the weights. Realizations can
therefore have different protocolizable populations.

Within each task-instance--policy--realization--model block, all arms receive
equal generation-request counts and randomized execution order under the same
frozen model and decoding policy. Repeated requests capture generation
variability, not additional independent tasks. Effect identification assumes
consistency, positive arm probabilities, and no cross-request interference.

\subsubsection{Independent Outcome Measurement}
\label{sec:method-measurement}
A prompt may produce valid but insecure code, code whose security cannot be
assessed, or no valid code at all. We keep these cases distinct. A security
\emph{Oracle} is the study's fixed, arm-blind procedure for measuring security
outcomes. The primary safety endpoint is \emph{oracle-evaluable secure-code yield}:
\begin{equation}
  Y=Y_CY_E\,\mathbb I\{\text{Oracle label is secure}\},
\end{equation}
where $Y_C$ indicates syntactically valid code and $Y_E$ indicates a supported
Oracle decision. The endpoint counts requests yielding valid code that the
security Oracle can evaluate and labels secure; it is not a security rate
conditioned on successful generation.

The Oracle must be qualified for the admitted languages,
task families, and code forms before formal use. Qualification includes
representative secure, insecure, and unknown examples and checks for plausible
differential errors. Separately, every syntax-valid program receives blinded
functional review against the complete original task, not the edited prompt.
Whole-task pass requires assessable material requirements and their fulfillment;
unresolved requirements cannot be credited as satisfied.

No-code, invalid-code, and Oracle-unknown outcomes contribute no observed secure
yield but keep their distinct labels. Unknown is not relabeled insecure.
Missing infrastructure evidence is not a zero outcome: it blocks the affected
analysis pending repair or replay of the same assignments. Every committed
assignment otherwise remains under its assigned arm. Fidelity, semantic
compliance, generation success, and non-target drift are diagnostics, not
filters that redefine the analyzed population.

Validity, Oracle coverage, unknown outcomes, functionality, and
secure-and-functional joint success accompany the primary endpoint. Coverage
sensitivity bounds distinguish gains in observed secure yield from changes in
unobserved security outcomes. Thus, a higher secure yield alone is not reported
as a lower underlying insecure-code rate when coverage changes.

\subsubsection{Policy-Effect Estimation}
The independent unit is the deduplicated task unit, not a generation request.
For a fixed hypothesis $h$ and generator model, let $\bar Y_{ia}$ first average
arm $a$'s request outcomes within each task instance and then combine instances
within task unit $i$ using their frozen weights. A realization $r\in R_h$ has
a prespecified weight $q_{hr}>0$, with $\sum_r q_{hr}=1$. If $n_{hr}$ eligible
task units are allocated to that realization, the arm mean is
\begin{equation}
  \widehat\mu_{ha}
  =\sum_{r\in R_h}q_{hr}
     \left(\frac{1}{n_{hr}}\sum_{i:r_{hi}=r}\bar Y_{ia}\right).
\end{equation}
We average within each realization before applying its original weight.
This preserves the declared policy mixture rather than letting realized sample
counts or failures choose new weights.

The Atomic assigned-arm intention-to-treat (ITT) effect is the risk difference
between target and no-op arm means:
\begin{equation}
  \widehat\tau_h=\widehat\mu_{hT}-\widehat\mu_{hN}.
\end{equation}
Unlike the observational RD used for prioritization, this contrast uses outcomes
from assigned prompt interventions on held-out tasks. For a Pair, the primary
RD interaction is the difference between two such contrasts:
\begin{equation}
  \widehat\delta_h
  =(\widehat\mu_{h,11}-\widehat\mu_{h,01})
   -(\widehat\mu_{h,10}-\widehat\mu_{h,00}).
\end{equation}
It asks whether the change associated with targeting one factor differs when
the other factor is also targeted. A nonzero interaction alone does not show
that the joint-target arm improves security. Both estimands remain specific
to the frozen policy, population mixture, and generator model; they are not
universal requirement effects or mediation claims.

Uncertainty is computed by resampling whole task units within original source
strata, carrying dependent arms, instances, effects, and models together and
preserving the frozen realization/source weights. Atomic primary contrasts
and Pair interactions use two separate simultaneous max-$|T|$ families.
Their intervals distinguish positive-meaningful, negative-meaningful,
practically null, and inconclusive effects relative to prespecified margins.
Failed provenance, inadequate support, excessive unknown coverage, or invalid
bootstrap support produces an explicit non-evaluable status, not an
unadjusted fallback. Secondary specificity and functionality tests require
separately frozen inference rules.

Stage II supplies a model-bound policy-effect record linked to the source
requirement, selected hypothesis, concrete edits, assignments, and measurements.
These links allow an explanation to state exactly what was compared and what
the resulting evidence supports. The observational reason for selection remains
separate from the randomized finding, regardless of its direction or evaluability.

\subsection{Stage III: Evidence-Grounded Explanations and Guidance}
\label{sec:method-guidance}

An effect estimate does not by itself tell a developer where to edit a prompt
or whether the change is appropriate. We specify a fixed explanation format
over the linked source, policy, and inference records. It connects the meaning
of a requirement to a tested modification and the action its evidence can
support. The explanation concerns observable effects of prompt policies,
not the generator's internal reasoning.

\subsubsection{Evidence-to-Explanation Mapping}
The explanation begins with the applicable task context and the requirement's
operation, object, and condition bindings from Stage I. It identifies the
source location or operation at which the requirement is expressed or would
be added, together with the non-target behavior to preserve. The concrete
target edit and its comparator then define the decision being explained.
For an Atomic policy this is target versus operation-matched no-op; for a
Pair, the explanation identifies both requirements and their four-cell design.
An illustrative edit cannot silently replace the frozen policy mixture whose
effect is reported.

The evidence portion distinguishes why the requirement was investigated from
what the intervention established. Structural admission and observational RD
provide the selection rationale. Randomized evidence provides the effect
estimate, simultaneous interval, practical margin, and effect status for the
tested generator and task population. Each statement draws on its corresponding
record: semantic relationships do not become causal edges, and an observational
rank does not become an estimated intervention effect. Missing evidence is
identified as unavailable rather than supplied by a plausible narrative.

The explanation reports functionality and measurement limits alongside the
security finding. It retains validity, Oracle coverage, unknown outcomes,
functional pass/fail/unknown assessments, and joint success without treating
them as interchangeable. Source incompleteness alone does not determine the
functional label; the independent review of generated code against the original
task does. Effect statements retain their model, endpoint, and frozen
realization/population mixture. A task-local binding makes the explanation
specific, but does not turn a population-average effect into a prediction for
that individual task or establish transfer to a new task population.

\subsubsection{Conditions for Security Guidance}
Guidance states what action the available evidence warrants, including when
no adoption advice can be given. Without formal randomized evidence, the
output is an unconfirmed modification option with a reason to investigate it.
Neither a high observational rank nor structural admission warrants describing
that option as validated. With evaluable randomized evidence, the explanation
can report the supported endpoint-specific effect even if other evidence
needed for adoption is missing.

A recommendation to adopt a tested edit additionally requires resolved
applicability, the exact tested comparison, and security and task-preservation
evidence meeting prospectively specified decision criteria. A favorable primary
effect alone is insufficient: an increase in oracle-evaluable secure-code yield
does not establish reduced underlying insecurity or preserved functionality.
Likewise, absence of a detected functional loss does not establish preservation.
The required secondary margins, coverage conditions, and joint decision rules
must be fixed before inspecting the relevant confirmation outcomes. Where a
required rule or result is unavailable, the explanation states that adoption
advice is withheld and identifies the missing support; it still reports the
available effect and uncertainty.

For an Atomic contrast, negative-meaningful evidence supports a warning about
the tested policy's effect on the named endpoint. Practically null evidence
places that average policy effect within the declared margin and scope.
For a Pair, these statuses instead describe the interaction: a negative or null
interaction does not imply that the joint edit is harmful or ineffective.
Inconclusive and non-evaluable evidence leave the corresponding estimand
unresolved for different reasons.
These findings are retained rather than converted into positive recommendations
or omitted. A harmful removal cannot validate the reverse addition policy.
For a Pair, a nonzero interaction describes how the two effects differ; advice
to apply both requirements needs separately supported joint-benefit and
task-preservation conclusions, not an interaction label alone.

In the SQL example, the resulting explanation links the supplied identifier
and deletion operation to a proposed explicit parameterization requirement,
states which deletion and return behavior the edit must preserve, and names
the matched no-op comparison. It then reports the evidence available for that
policy and its limits. Before confirmation, this is a task-grounded option to
test. After confirmation, the same explanation format can present a supported
effect or a reason to withhold advice. It never infers that the original code
was insecure merely because the prompt left parameterization unexpressed.

This reporting step does not revise scientific effect statuses, remove
assigned tasks, or rerank hypotheses. Computational evaluation assesses the
evidence underlying proposed changes; the separate expert study assesses the
perceived quality and utility of the explanation products.
